\chapter {DQN implementation}

The complete source code reside in github location: \url{https://github.com/samiislam/LearningToPlayPong/tree/main/code/dqn}

\section{The DQN Model}
\lstinputlisting[language=Python, caption={code\textbackslash dqn\textbackslash core\textbackslash model.py}, label={code:dqn_core_model}]{code/dqn/core/model.py}

\section{Gymnasium Wrappers}
\lstinputlisting[language=Python, caption={code\textbackslash dqn\textbackslash core\textbackslash wrappers.py}, label={code:dqn_core_wrappers}]{code/dqn/core/wrappers.py}

\section{The Agent}
\lstinputlisting[language=Python, caption={code\textbackslash dqn\textbackslash core\textbackslash dqn\_agent.py}, label={code:dqn_core_dqn_agent}]{code/dqn/core/dqn_agent.py}

\section{Training Pong}
\lstinputlisting[language=Python, caption={code\textbackslash dqn\textbackslash dqn\_pong.py}, label={code:dqn_dqn_pong}]{code/dqn/dqn_pong.py}

\section{Playing Pong}
\lstinputlisting[language=Python, caption={code\textbackslash dqn\textbackslash dqn\_play.py}, label={code:dqn_dqn_play}]{code/dqn/dqn_play.py}

\section{Engineering and Training Speedup Optimizations}
\label{sec:dqn-optimizations}
Six engineering and speedup code optimizations in $\hyperref[code:dqn_dqn_pong]{dqn\_pong.py}$:
\begin{itemize}
\item Vectorized environments [line 83], where N\_ENVS=8 number of parallel pong environments are used to collect 8 transitions per iteration in parallel, increasing the environment throughput without additional GPU cost. 
\item Polyak averaging, also known as soft target updates, which continuously blends target weights toward the online network, avoiding the periodic stale-target intervals that slow early learning [lines 187-188]. 
\item Mixed-precision training that halves memory bandwidth and doubles GPU arithmetic throughput by running most operations in fp16 instead of fp32 [line 57, 121]. 
\item Compiling the pytorch environment using cudagraphs [lines 99, 107; $\hyperref[code:dqn_core_dqn_agent]{dqn\_agent.py}$ line 99] which captures the entire forward pass as a single CUDA graph, eliminating per-step CPU kernel-launch overhead. 
\item Non-blocking tensor transfer [$\hyperref[code:dqn_core_dqn_agent]{dqn\_agent.py}$ line 150] which overlaps CPU-to-GPU data transfers with computation, reducing idle time on both sides. 
\item A different set of hyperparameters [lines 29-42].
\end{itemize}